\documentclass{article}

\usepackage{amsmath,amssymb}
\usepackage{xurl}
\usepackage[hidelinks]{hyperref}

% Shared commands for notes in docs/paper-gaps/.

\DeclareUnicodeCharacter{2080}{\ifmmode{}_{0}\else\textsubscript{0}\fi}
\DeclareUnicodeCharacter{2081}{\ifmmode{}_{1}\else\textsubscript{1}\fi}
\DeclareUnicodeCharacter{2082}{\ifmmode{}_{2}\else\textsubscript{2}\fi}
\DeclareUnicodeCharacter{2099}{\ifmmode{}_{n}\else\textsubscript{n}\fi}

\newcommand{\Lean}{\textsc{Lean}}
\ExplSyntaxOn
\NewDocumentCommand{\leanid}{m}{
  \ifmmode
    \textsf{\detokenize{#1}}
  \else
    \group_begin:
    \urlstyle{sf}
    \tl_set:Nn \l_tmpa_tl {#1}
    \tl_replace_all:Nnn \l_tmpa_tl {\_} {_}
    \exp_args:NV \path \l_tmpa_tl
    \group_end:
  \fi
}
\ExplSyntaxOff
\providecommand{\tr}{\operatorname{tr}}
\newcommand{\tnleanIssueBase}{https://github.com/LionSR/TNLean/issues/}
\newcommand{\tnleanPrBase}{https://github.com/LionSR/TNLean/pull/}
\newcommand{\tnleanDocBase}{https://sirui-lu.com/TNLean/docs/}
\newcommand{\ghissue}[1]{\href{\tnleanIssueBase#1}{GitHub issue \##1}}
\newcommand{\ghpr}[1]{\href{\tnleanPrBase#1}{PR \##1}}
\newcommand{\doclink}[1]{\href{\tnleanDocBase#1.html}{\path{#1}}}

% Dirac notation and the identity operator, as in blueprint/src/macros/common.tex.
% \providecommand keeps note-local definitions authoritative.
\usepackage{dsfont}
\providecommand{\ket}[1]{|#1\rangle}
\providecommand{\bra}[1]{\langle #1|}
\providecommand{\braket}[2]{\langle #1 | #2 \rangle}
\providecommand{\ketbra}[2]{|#1\rangle\!\langle #2|}
\providecommand{\Id}{\mathds{1}}
\providecommand{\one}{\mathds{1}}
\providecommand{\C}{\mathbb{C}}
\providecommand{\MN}[1]{M_{#1}(\C)}
\providecommand{\MD}{M_D(\C)}

% External referencing for paper sources.  The build helper writes sanitized
% auxiliary files under build/paper-aux/ and excludes duplicated source labels.
\usepackage{xr}

% Import a generated paper auxiliary file with a caller-supplied label prefix.
% The candidate paths support builds from docs/paper-gaps/, the repository root,
% and build/paper-aux/ itself.
\newcommand{\importpaperaux}[2]{%
  \IfFileExists{../../build/paper-aux/#2.aux}{%
    \externaldocument[#1]{../../build/paper-aux/#2}%
  }{%
    \IfFileExists{build/paper-aux/#2.aux}{%
      \externaldocument[#1]{build/paper-aux/#2}%
    }{%
      \IfFileExists{#2.aux}{%
        \externaldocument[#1]{#2}%
      }{}%
    }%
  }%
}

% General external reference macros
\newcommand{\paperref}[2]{\ref{#1-#2}}
\newcommand{\papereqref}[2]{(\ref{#1-#2})}

% CPSV16 source (1606.00608 / MPDO-22-12-17-2.tex) external document and shortcuts
\importpaperaux{cpsv16-}{1606.00608/MPDO-22-12-17-2-tnlean}
\newcommand{\cpsvref}[1]{\paperref{cpsv16}{#1}}
\newcommand{\cpsveqref}[1]{\papereqref{cpsv16}{#1}}

% Verdict marker: \gapnote{<kind>}{<status>}. Kind is the policy
% classification (clarification, local-correction, scope-restriction,
% unfaithful, false-source, open-gap); status is open, wip (elimination
% underway), resolved, or historical. Machine-read by the site index;
% prints a verdict line.
\newcommand{\gapnote}[2]{\par\noindent\textbf{Verdict:} #1 (#2).\par}


\title{The Implication Label in the Proof of MPDO Theorem 4.9}
\author{}
\date{2026-07-12}

\begin{document}
\maketitle

\section{The cyclic mismatch}

The statement of Theorem~4.9 in
\cite[lines~851--893]{Cirac2016MPDO_arXiv} numbers SAL together with literal
ZCL as condition~(ii), Gibbs form together with literal ZCL as condition~(iii),
the BNT physical-sector structure as condition~(iv), and the
renormalization fixed-point property as condition~(v). The propositions in
the Appendix have the following literal component conclusions:
\begin{align*}
  \mathtt{prop2to3}&:\quad (ii)\Longrightarrow(iv),\\
  \mathtt{prop3to4}&:\quad (iv)\Longrightarrow
      \text{the GSNNCH-form clause of (iii)},\\
  \mathtt{prop4to2}&:\quad \text{the GSNNCH-form clause of (iii)}
      \Longrightarrow\text{the SAL clause of (ii)},\\
  \mathtt{prop2to5}&:\quad (ii)\Longrightarrow
      \text{a channel pair between the two- and three-site tensors}.
\end{align*}
These are the displayed statements at source lines~1740--1817. In particular,
\texttt{prop3to4} does not prove the ZCL conjunct of condition~(iii), and
\texttt{prop4to2} does not prove that conjunct of condition~(ii); the ZCL
statement is carried from the surrounding implication. Likewise,
\texttt{prop2to5} does not itself conclude condition~(v). The concluding proof
at lines~1819--1825 nevertheless assigns the propositions successively to
$(ii)\Rightarrow(iii)$, $(iii)\Rightarrow(iv)$,
$(iv)\Rightarrow(ii)$, and $(iii)\Rightarrow(v)$. Thus the discrepancy concerns
the proof path and the component conclusions, and is not repaired by changing
only the last occurrence from ``(iii)'' to ``(iv)''.

Proposition~\texttt{3to5} at lines~1510--1563 constructs a pair of channels
from one supplied physical-sector factorization.  Proposition
\texttt{prop2to5} then says that the same construction is first performed in
each locally orthogonal BNT subspace. Its assertion that the corresponding
maps can be combined is the additional all-BNT step. Lines~1821--1825 then
apply the two-to-three-site channels twice to obtain the one-to-two-site maps
for the blocked tensor in condition~(v).

\section{Formal treatment}

The formal development therefore does not attach the one-factorization
construction to any of the incorrect labels in the concluding paragraph.  It
cites Proposition~\texttt{3to5} and the twice-applied channel observation at
lines~1821--1825 directly.  The resulting theorem is conditional on one
supplied factorization and is not the full all-BNT implication of
Theorem~4.9.  The latter still requires a direct-sum construction controlled
by the mutually orthogonal outer physical supports.

Files that use the concluding paragraph carry a \emph{Local fix (cyclic
implication labels)} marker and cite this note. This records proof-path drift
in the printed proof without altering any mathematical
hypothesis.  The separate restriction to one supplied factorization is
recorded in
\path{docs/paper-gaps/cpgsv17_pf_rank_one.tex} and tracked in issue~\#6632.

\bibliographystyle{alpha}
\bibliography{references}

\end{document}
